\documentclass[12pt,a4paper,oneside]{article} % part, *section, *paragraph
%\documentclass[12pt,a4paper,oneside]{report} % part, *chapter, *section, *paragraph
%\documentclass[12pt,a4paper,oneside]{book} % part, *chapter, *section, *paragraph

\usepackage[english]{babel}
\usepackage[latin2]{inputenc}
\usepackage{setspace}
\usepackage{float}
\usepackage[T1]{fontenc} %\showhyphens{integetnem} \showhyphens{meggy�ltetv�ny}
%\usepackage{ucs}\usepackage[utf8x]{inputenc}
%\usepackage{fullpage} %\usepackage[margin=2cm]{geometry}
%\usepackage{amsmath,amssymb}
\usepackage{url} %VAGY %\usepackage[unicode]{hyperref}
\usepackage{listings}
%\usepackage{color} 
\usepackage{graphicx}
\usepackage{times}
\usepackage{fullpage}
%\usepackage{showkeys} % \label{} ki�r�sa \ref{}hez
 
\let\stdsection\section
\renewcommand\section{\clearpage\stdsection}
\begin{document}
\onehalfspacing
\begin{center}
\vspace{2em}
%\large{\bf{}BUDAPESTI M\H{U}SZAKI \'ES GAZDAS\'AGTUDOM\'ANYI EGYETEM}
\vspace{.5em}

\includegraphics[width=60mm]{10000000000004C00000014A5D3CE278.eps}

Budapest University of Technology and Economics

Department of Measurement and Information Systems
%\vspace{.5em}

%\large{\bf{}VILLAMOSM\'ERN\"OKI \'ES INFORMATIKAI KAR\\
%M\'ERN\"OK INFORMATIKUS SZAK}

\vspace{2em}

%\Large{Informatikai Technol\'ogi\'ak szakir\'any\\
%Rendszerfejleszt\'es \'agazat}

%\vspace{1.5em}

%\Large{\bf{}\"On\'all\'o labor (BMEVIIIA337)}

%\vspace{1.5em}

\LARGE{\bf{}RTF tokenizer for LibreOffice writerfilter (rewrite of the RTF import)}

\vspace{2em}

\begin{tabular}{ l c}
  %hopefully i can avoid the photo this time
  %\parbox{30mm}{\includegraphics[width=20mm]{20090926.eps}} &
  \parbox{120mm}{
\begin{center}
\normalsize{Vajna Mikl�s (AYU9RZ),

4th semester, (MSc) computer engineering student

Consultants: C�dric Bosdonnat, M�t� Andr�s, Novell

Internal consultant: Horv�th �kos, MIT

Specialization in Dependable System Design

Internship Report

2011/12. 1st semester}
\end{center}
} \\
\end{tabular}


\end{center}
\thispagestyle{empty}
\newpage

\tableofcontents
\newpage

\section{Introduction}

LibreOffice is a power-packed free, libre and open source personal
productivity suite for Windows, Macintosh and GNU/Linux, that gives you six
feature-rich applications for all your document production and data processing
needs: Writer, Calc, Impress, Draw, Math and Base\cite{libreoffice}.

During my internship I worked on Writer. Writer has support for lossless import
and export from/to the ODF file format (.odt). Additionally several other
filters are supported which are marked as "alien": such filters may lose some
information upon loading/saving. Such an alien RTF\cite{rtf} filter was already
implemented in LibreOffice 3.4, and my task was to improve it.

The source code of LibreOffice is divided to several (at the time of writing:
225) modules, one of them is writerfilter, which contains the import filter for
the DOCX file format.

The basic idea was that the writerfilter module already contains code to map
Microsoft Office's concepts to LibreOffice's concepts (for example Microsoft
Office uses section breaks for different page settings, LibreOffice uses page
styles for that purpose, etc.), so in case I can implement a new RTF filter
within writerfilter, we can avoid code duplication and gain mappings which were
missing in the old RTF filter.

The rest of this paper is structured as follows. First, I introduce necessary
background knowledge, which was present before the current paper, but is needed
to understand the rest of this work (Section 2). Section 3 describes the design
of the solution, as well as relation with the underlying techniques. Next, I
detail the implementation I created (Section 4) and also evaluate it (Section
5). Finally I give a summary, including future development directions (Section
6).

\section{Background}

A filter in LibreOffice is a class, that implements given interfaces, and
performs the import or the export (or both) of a file format to one of the six
applications.

Registrations of such classes are usually handled via UNO (Universal Network
Objects), the interface based component model of LibreOffice.\footnote{Old
filters are not using UNO and are bundled together with applications (e.g.
Writer).} Each filter has a configuration file, which lists the file format
detector service and the filter service.

I did not have to pay attention to the RTF file format detector service, as it
was already working and I had no plans to improve it. The filter was an UNO
one, but the import part of it was just a stub, that called the old RTF
importer, that was a builtin one.

Implementing a filter via UNO has multiple advantages:

\begin{itemize}
\item The applications are written in C++, but the implementation of UNO services can be anything supported by UNO (e.g. Java or Python).
\item The filter can be built once the UNO interfaces implemented by the
application are available, even if the source code of the application it not
available\footnote{This may be strange as a goal for a free software project,
but given that LibreOffice is huge and building it is non-trivial, it's
useful.}.
\end{itemize}

The already cited RTF specification defines the major tasks of an RTF reader:

\begin{itemize}
\item Separate text from RTF controls.
\item Parse an RTF control.
\item Dispatch an RTF control.
\end{itemize}

Separating text from RTF controls is relatively simple: all RTF controls begin
with a backslash and can contain a defined set of characters only.

Parsing an RTF control is also relatively simple: RTF defines different type of
control words and each of them a simple rule of what type of parameters can
they take (if any).

Dispatching an RTF control, on the other hand, is relatively complicated: once
needs having a really good knowledge of the API of the target application,
since most of the RTF control words will result in an API call, passing the
parameters of the control word to the document model.

Next to control words, RTF has the concept of groups. They start with a
\textbraceleft and end with a \textbraceright -- and they should be threated
like a stack: the scope of control words starts where they appear and ends at
the end of the current group.

RTF defines the following control word types:

\begin{itemize}
\item flags: they take no parameter, for example \textbackslash{}sbknone
\item destinations: these appear at the start of a group and if they are
	ignored, the whole group should be ignored, for example \textbackslash{}fonttbl
\item symbols: they don't take parameters, either -- special symbols like \textbackslash{}tab have this type
\item toggles: they switch a property on or off, for example bold characters are marked using \textbackslash{}b
\item value: they always take a parameter, for example font size (\textbackslash{}fs)
\end{itemize}

Now that we are aware of basic concepts of RTF, the next section overviews the
API requirements of Writer and explains the design of the filter that can parse
and RTF document.

\section{Design}

Bla bla.

\section{Implementation}

Bla bla.

\section{Testing}

Bla bla.

\section{Future work}

Bla bla.

\clearpage

\begin{thebibliography}{4}
\addcontentsline{toc}{section}{References}

\bibitem{libreoffice} Libreoffice, http://libreoffice.org/

\bibitem{rtf} Word 2007: Rich Text Format (RTF) Specification, version 1.9.1, http://bit.ly/qD1CTV

\end{thebibliography}

\end{document}
